\clearpage
\subsection{Learning rate: follow one parameter}
A parameter is a trainable weight or bias. The vector $\theta$ collects all of
them. A batch is a subset of training samples; its loss $C_B$ is their average
prediction error. The gradient points locally uphill in parameter space, so an
optimizer takes a step in the opposite direction. The update index $k$ counts
parameter changes, not epochs and not physical CFD time steps.
\[
g_k=\nabla_\theta C_B(\theta_{k-1}),\qquad
\theta_k=\theta_{k-1}-\eta g_k.
\]
The learning rate $\eta>0$ sets the step scale. For one scalar parameter, take
\[
C(\theta)=\tfrac12(\theta-3)^2,\qquad
g=\frac{\mathrm dC}{\mathrm d\theta}=\theta-3.
\]
Starting from $\theta_0=1$ gives $g_1=-2$. With $\eta=0.5$, the next parameter
is $\theta_1=1-0.5(-2)=2$: the negative gradient moves it toward the minimum
at $\theta=3$. Repeating the update gives
\[
\underbrace{\theta_k-3}_{e_k}
=(1-\eta)\underbrace{(\theta_{k-1}-3)}_{e_{k-1}}.
\]
Thus this quadratic gradient-descent example converges when $|1-\eta|<1$,
or $0<\eta<2$. For $\eta=0.05$ it approaches slowly; for $\eta=0.5$ its error
halves each update; for $\eta=2.1$ its error changes sign and grows.
This bound depends on this loss's curvature. It is not a general neural-network
or Adam learning-rate limit.

\textbf{Notebook experiment.} Plot the parameter and the loss against update
number for these rates, including $\eta=1.5$ (alternating, decaying errors). Keep the starting value and loss fixed so that
the effect of the rate can be identified. The experiment uses NumPy and does
not retrain the fluid surrogate.

\subsection{Adam: two memories for every parameter}
Adam stores a running average $m_k$ of signed gradients and a running average
$v_k$ of squared gradients. Both arrays have the same shape as $\theta$.
The symbol $v_k$ in this optimizer is memory, not a fluid velocity.
Starting with $m_0=v_0=0$,
\begin{align*}
m_k&=\beta_1m_{k-1}+(1-\beta_1)g_k,\\
v_k&=\beta_2v_{k-1}+(1-\beta_2)g_k^2.
\end{align*}
The parameters $0\leq\beta_1,\beta_2<1$ control memory. The notebook uses
the Keras defaults $\beta_1=0.9$ and $\beta_2=0.999$. Multiplication, squaring,
division and square roots operate separately on every component.
The first memory smooths the direction; the second measures recent gradient
magnitude. Dividing by a square-root magnitude makes the step scale adaptive
for each parameter. The second memory averages squared gradients; it is not a statistical variance. It does not remove the base learning-rate choice.

\clearpage
\subsection{Bias correction and the implemented update}
The zero initial memories underestimate the early average magnitudes. After
$k$ updates their bias-corrected versions are
\[
\widehat m_k=\frac{m_k}{1-\beta_1^k},\qquad
\widehat v_k=\frac{v_k}{1-\beta_2^k}.
\]
The familiar Adam algorithm writes
\[
\theta_k=\theta_{k-1}-\eta\frac{\widehat m_k}
{\sqrt{\widehat v_k}+\epsilon_{\rm alg}}.
\]
Here the small positive $\epsilon_{\rm alg}$ prevents division by zero.
For exact reproduction, Keras uses the paper's \emph{epsilon-hat} convention:
\[
\alpha_k=\eta\frac{\sqrt{1-\beta_2^k}}{1-\beta_1^k},\qquad
\boxed{\theta_k=\theta_{k-1}-\alpha_k
\frac{m_k}{\sqrt{v_k}+\epsilon_K}.}
\]
Its default is $\epsilon_K=10^{-7}$. The equivalent denominator in the
bias-corrected form contains $\epsilon_K/\sqrt{1-\beta_2^k}$; the two epsilon
symbols should not be assigned the same value when reproducing an exact step.

\textbf{One numerical Adam step.} For the preceding quadratic, $\theta_0=1$
and $g_1=-2$. Hence
\[
m_1=0.1(-2)=-0.2,\qquad v_1=0.001(4)=0.004,
\]
\[
\widehat m_1=-0.2/0.1=-2,\qquad
\widehat v_1=0.004/0.001=4.
\]
Use $\eta=0.1$ only for this visible arithmetic example. Substitution into the
Keras-form update gives $\theta_1\simeq1.09999984$. The notebook executes
this calculation and checks its value. The fluid surrogate keeps its separate
learning rate of $0.001$.

Using the same $\eta=0.1$, gradient descent would move by $0.2$, giving
$\theta_1=1.2$, while Adam moves by approximately $0.1$. With uncorrected
memories, the alternative update $-\eta m_1/(\sqrt{v_1}+\epsilon_K)$ would
move by approximately $0.316$. Bias correction changes the early step scale.

\clearpage
\subsection{Read the actual training settings}
The 5000 samples split into 3500 training, 750 validation and 750 test samples.
The $4$--$32$--$32$--$1$ network has
$4(32)+32+32(32)+32+32(1)+1=1249$ trainable parameters.
Each epoch contains 55 updates: 54 batches of 64 samples and a final batch of 44.
Running all 500 epochs would give 27,500 updates.

\begin{center}\begin{tabular}{p{.31\linewidth}p{.61\linewidth}}\toprule
Setting & Meaning in the Week 2 training call\\\midrule
\texttt{learning\_rate=1e-3} & Base Adam step scale $\eta=0.001$.\\
\texttt{batch\_size=64} & Up to 64 samples contribute to each gradient update.\\
One epoch & One pass through the training set: approximately $\lceil N_{\rm train}/64\rceil$ updates, including a smaller final batch.\\
\texttt{epochs=500} & Maximum number of passes, not a promise to run all 500.\\
\texttt{patience=30} & Stop after 30 epochs without lower validation loss (default improvement threshold: zero).\\
\texttt{restore\_best\_weights=True} & Restore weights from the best monitored epoch.\\\bottomrule
\end{tabular}\end{center}
If epoch 120 has the lowest validation loss and the next 30 epochs do not improve it, training ends after epoch 150 and restores epoch-120 weights. Early stopping does not reduce the learning rate. For each independent rate trial, call \texttt{keras.utils.set\_random\_seed(7)}
before rebuilding the model, optimizer and callback. Hold the split,
training-only scaling, architecture and batch policy fixed; compare matched seeds. Select using validation loss; reserve test data for final evaluation.
Training loss alone does not establish physical accuracy or extrapolation.

\textbf{Sources.} Kingma and Ba, \emph{Adam: A Method for Stochastic Optimization}
(ICLR 2015), \url{https://arxiv.org/abs/1412.6980}; Keras Adam documentation,
\url{https://keras.io/api/optimizers/adam/}.
\clearpage
